\begin{definition} \label{definition:derivation_in_an_algebra_over_a_ring}
    Let $R$ be a \CrefAndHyperrefIfExist{definition:ring}{(not-necessarily commutative) ring with unity} and let $A$ be an \CrefAndHyperrefIfExist{definition:algebra_of_a_ring}{$R$-algebra}. Let $M$ be an \CrefAndHyperrefIfExist{definition:bimodule_over_rings_module_over_a_ring}{$A$-module}. 
    A \hldef{derivation on $A$ over $R$ with values in $M$} is an $R$-linear map $D: A \to M$ satisfying the following properties for all $a, b \in A$:
    \begin{enumerate}
        \item Additivity: $D(a + b) = D(a) + D(b)$,
        \item Leibniz rule: $D(ab) = aD(b) + bD(a)$.
    \end{enumerate}

    If $M = A$, then $D$ is called a \hldef{derivation of $A$}. The set of all derivations from $A$ to $M$ is denoted by \hl{$\mathrm{Der}_R(A, M)$}, which naturally carries the structure of an $A$-module.
\end{definition}
